%% =====================================================================
%%  T-09-01  A Reputation You Built on Purpose
%%  -------------------------------------------------------------------
%%  One idea per file. Core is mandatory. Slots in canonical order.
%%  Max 700 words. No layout commands. No filler. New source material
%%  for this entry goes in sources/<BOOK>/T-09-01.tex -- never by
%%  rewriting this file (docs/RULES.md R0.1).
%%  =====================================================================
%% @kind: topic
%% @id: T-09-01
%% @title: A Reputation You Built on Purpose
%% @subtitle: Others cannot see you; they see the standing they assigned you
%% @chapter: c09
%% @order: 10
%% @status: stable
%% @sources: B3
%% @updated: 2026-09-19

\topic{T-09-01}{A Reputation You Built on Purpose}{Others cannot see you; they see the standing they assigned you}

\begin{core}
People cannot inspect character directly, so they price you by the standing they already hold in memory and read each new deed through it. The same action reads as brilliance from one name and recklessness from another. A reputation built deliberately --- chosen traits, displayed consistently --- is therefore not vanity but leverage: it answers the judgment before you enter the room.
\end{core}
\begin{mechanism}
Judgments run on appearances --- clothes, gestures, words, visible results --- because inner character is unfathomable and first-hand checking is tiring. A settled reputation spares that work, which is why it is believed. It also shields: attention fixes on the published image, and scrutiny that would otherwise reach the person is absorbed by the file. The same mechanism is the attack surface. Seeded doubt puts its target in a dilemma with no good move: deny it and the denial itself reads as suspicious --- why defend so hard? --- while ignoring it leaves the doubt standing unrefuted. An attacker who already holds standing of their own can escalate from doubt to light ridicule, which signals security, puts the target permanently on the defensive, and draws attention to the mocker.
\end{mechanism}
\begin{tells}
  \item Praise or distrust that arrives ahead of any dealing --- people are reacting to the file, not the person.
  \item An image maintained with a consistency too careful to be accidental.
  \item Under a rumour: either escalating self-defence or studied silence. Both responses pay the rumour.
  \item Ridicule of a rival from someone whose own standing is already secure.
  \item Standing quoted as its own evidence, with no visible works behind it.
\end{tells}
\begin{moves}
  \item Pick two or three traits and make every visible act consistent with them. A reputation is built by public repetition, not by assertion.
  \item Treat the first doubt as an event: answer it once, early, with evidence, then stop. The denial spiral is the trap.
  \item Never let the image be assigned by default. Silence lets others write the file.
  \item Against a rival, sow specific doubt first; move to ridicule only from standing you already hold.
  \item Keep part of the real self out of view. The published image should be the least interesting thing about you.
\end{moves}
\begin{counters}
  \item Price deeds, not names: re-read the work as if the name were absent, and see what survives.
  \item If you catch yourself defending against a rumour for the second time, you are paying it. One evidence-backed statement, then be seen doing the work.
  \item Watch for consensus about a person that everyone repeats and nobody checked --- manufactured agreement is covered in T-13-01, and a seeded room is covered in T-23-01.
  \item Notice borrowed symbols doing the pricing --- titles, offices, rooms. Ask what remains of the standing without them (T-18-02).
\end{counters}
\begin{cost}
The image and the person diverge, and the gap must be patrolled forever: every visible act has to fit the file, a permanent tax on behaviour. Because the reputation is the asset, exposure as manufactured re-prices everything previously credited to it --- the collapse takes the past with it. And the mechanism is exactly the one con artists run on: judging by standing instead of works is what a manufactured reputation exists to exploit.
\end{cost}
\begin{field}
Wherever strangers price strangers: hiring, court politics, public markets for trust. Most total in small communities, where the file is shared infrastructure and one rumour re-prices the system.
\end{field}

\sources{B3}
\seesources
\seealso{T-13-01, T-18-02, T-23-01}
